\begin{mistake}{2026-07-15}{03}

\begin{problemcontent}
设立体图形的底是介于 \(y=x^2-1\) 和 \(y=0\) 之间的平面区域，而它的垂直于 \(x\) 轴的任一截面是等边三角形，求立体体积 \(V\)。
\end{problemcontent}

\begin{solutioncontent}
底区域由 \(y=x^2-1\) 与 \(y=0\) 围成。二者交点满足
\[
x^2-1=0,
\]
故 \(x=\pm1\)，积分区间为 \([-1,1]\)。

截面垂直于 \(x\) 轴，因此固定 \(x\) 后，截面等边三角形的边长等于底区域在该 \(x\) 处的纵向长度：
\[
s=0-(x^2-1)=1-x^2.
\]

等边三角形面积公式为
\[
A=\frac{\sqrt{3}}{4}s^2,
\]
所以截面面积为
\[
A(x)=\frac{\sqrt{3}}{4}(1-x^2)^2.
\]

由截面法求体积：
\[
V=\int_{-1}^{1}A(x)\,dx
=\frac{\sqrt{3}}{4}\int_{-1}^{1}(1-x^2)^2\,dx.
\]
展开并利用偶函数性质：
\[
\int_{-1}^{1}(1-x^2)^2\,dx
=2\int_0^1(1-2x^2+x^4)\,dx.
\]
计算得
\[
2\left[x-\frac{2}{3}x^3+\frac{1}{5}x^5\right]_0^1
=2\left(1-\frac{2}{3}+\frac{1}{5}\right)
=\frac{16}{15}.
\]
因此
\[
V=\frac{\sqrt{3}}{4}\cdot\frac{16}{15}=\frac{4\sqrt{3}}{15}.
\]
\end{solutioncontent}

\mistakeknowledge{
\begin{itemize}
  \item \textbf{核心公式/定理}：截面法体积 \(V=\int_a^b A(x)\,dx\)；等边三角形面积 \(A=\frac{\sqrt{3}}{4}s^2\)。
  \item \textbf{易错点}：截面垂直于 \(x\) 轴时应固定 \(x\)，边长取底区域纵向长度 \(0-(x^2-1)=1-x^2\)，不是横向长度。
  \item \textbf{关联知识}：积分区间由两曲线交点确定；本题由 \(x^2-1=0\) 得 \(x\in[-1,1]\)。
\end{itemize}}

\end{mistake}
